\documentclass[11pt]{article}

\newcommand{\cnum}{Math 132H}
\newcommand{\ced}{Winter 2026}
\newcommand{\ctitle}[4]{\title{\vspace{-0.5in}\cnum, \ced\\Assignment #1: #2}\author{\vspace{-0.35in}\\#3, #4}}
\usepackage{enumitem}
\usepackage[usenames,dvipsnames,svgnames,table,hyperref]{xcolor}
\usepackage{amsmath, amsfonts}
\usepackage{graphicx} % Required for inserting images
\usepackage{float}
\usepackage{listings}
\usepackage{xcolor}
\usepackage{hyperref}
\usepackage{comment}

\renewcommand*{\theenumi}{\alph{enumi}}
\renewcommand*\labelenumi{(\theenumi)}
\renewcommand*{\theenumii}{\roman{enumii}}
\renewcommand*\labelenumii{\theenumii.}

\definecolor{codegreen}{rgb}{0,0.6,0}
\definecolor{codegray}{rgb}{0.5,0.5,0.5}
\definecolor{codepurple}{rgb}{0.58,0,0.82}
\definecolor{backcolour}{rgb}{0.95,0.95,0.92}
\definecolor{header}{HTML}{205459}
\definecolor{solution}{HTML}{204d52}

\newcommand{\solution}[1]{{{\textcolor{header}{Solution:} \textcolor{solution}{#1}}}}
\setlist[enumerate]{label=(\alph*)}

\lstdefinestyle{mystyle}{
    backgroundcolor=\color{backcolour},
    commentstyle=\color{codegreen},
    keywordstyle=\color{magenta},
    numberstyle=\tiny\color{codegray},
    stringstyle=\color{codepurple},
    basicstyle=\ttfamily\footnotesize,
    breakatwhitespace=false,
    breaklines=true,
    captionpos=b,
    keepspaces=true,
    numbers=left,
    numbersep=5pt,
    showspaces=false,
    showstringspaces=false,
    showtabs=false,
    tabsize=2
}


\begin{document}
\ctitle{\#2}{}{Asher Christian}{006-150-286}
\date{}
\maketitle

\section{Exercise 1.14}
\emph{
    Suppose $\{a_n\}_{n=1}^{N}$ and $\{b_n\}_{n=1}^{N}$ are two finite sequences of complex numbers.
    Let $B_k = \sum_{n=1}^{k}b_n$ denote the partial sums of the series $\sum_{}^{}b_n$ with $B_0 = 0$.
    Prove the summation by parts formula
    \[
        \sum_{n=M}^{N}a_nb_n = a_NB_N - a_MB_{M-1} - \sum_{n=M}^{N-1}(a_{n+1}-a_n)B_n
    \] 
}
\solution{
    By letting $c_n = a_n - a_{n-1}$ and  $C_n = a_n = \sum_{k=1}^{n}c_n$ we can rewrite what we are trying to prove as
    \[
        \sum_{n=M}^{N}C_nb_n = C_NB_N - C_MB_{M-1} - \sum_{n=m}^{N-1}c_{n+1}B_n
    \] 
    and we can assume without loss of generality that $a_0 = 0$ or define it to be so.
    expanding sums we get
    \[
        \sum_{n=M}^{N}\sum_{k=1}^{n}c_kb_n = \sum_{k=1}^{N}\sum_{j=1}^{N}c_kb_j - \sum_{k=1}^{M}\sum_{j=1}^{M-1}c_kb_j - \sum_{n=M}^{N-1}\sum_{k=1}^{n}c_{n+1}b_k
    \] 
    and moving terms so that everything is positive we now try to prove
    \[
        \sum_{k=1}^{N}\sum_{j=1}^{N}c_kb_j = \underbrace{\sum_{n=M}^{N}\sum_{k=1}^{n}b_nc_k}_{(a)} + \underbrace{\sum_{k=1}^{M}\sum_{j=1}^{M-1}c_kb_j}_{(b)} + \underbrace{\sum_{n=M}^{N-1}\sum_{k=1}^{n}c_{n+1}b_k}_{(c)}
    \] 
    Now, by drawing a picture it is clear that these sums are the same, however I will verify it here.
    Let $(i,j)$ be an index representing $c_i b_j$ such that  $1 \le i,j \le N$. Then there are three cases. If  $i \le M, j \le M-1$ then the term
    is in $(b)$
    and it is not in the other two sums since $j < M $ so it is not in  $(a)$ and $i < M+1$ so it is not in  $(c)$.
    If $j \ge M$ and  $i \le j$ then  $c_ib_j \in (a)$ by inspection and it is clearly not in  $(b)$ since  $j > M-1$ and it is not in  $(c)$ since $i \le j$.
    Lastly if $i \ge M+1$ and  $j < i$ then  $c_ib_j $ is  $(c)$ by similar reasoning as before and likewise the terms arent in  $(a)$ or  $(b)$ so we conclude that the 
    three different partial sums partition the set of all possible indices and so the left and right sides are the same.
}

\section{Exercise 1.16}
\emph{
    Determine the radius of convergence of the series $\sum_{n=1}^{\infty}a_nz^{n}$ when:
    \begin{enumerate}
        \item $a_n = (\log n)^2$
        \item $a_n = n!$
        \item  $a_n = \frac{n^2}{4^{n} + 3n}$ 
        \item $a_n = \frac{(n!)^{3}}{(3n)!}$
    \end{enumerate}
}
\solution{
    \begin{enumerate}
        \item We can use Hadamard's formula on $a_n = (\log n)^2$ we have $\frac{1}{R} = \limsup_{n\to \infty} (|a_n|)^{\frac{1}{n}} = \limsup_{n\to \infty}(\log n)^{\frac{2}{n}} = 1$ 
            Since $\log n > 1$ for large  $n$ and the power is positive so it is greater than or equal to 1 and since  $(\log n)^2 < n$ for large $n$ we get it is less than or equal to $1$ so it is  $1$ and so  $ R = 1$
        \item Again we consider
            \[
            \frac{1}{R} = \limsup_{n\to \infty} (n!)^{\frac{1}{n}} = +\infty
            \] 
            so the radius of convergence for this series is 0
        \item 
            For this one we recognize that
            \[
            \limsup_{n\to \infty}\left(\frac{n^2}{4^{n}+3n}\right)^{\frac{1}{n}} = \frac{1}{4}
            \] 
            and taking reciprocal we get $R = 4$
        \item using the approximation we get
            \begin{align*}
                \frac{ (n!)^{3}}{(3n)!} &\approx \frac{(nc^{n+\frac{1}{2}}e^{-n})^{3}}{c(3n)^{3n+\frac{1}{2}}e^{-3n}}\\
                                        &= \frac{c^{3}e^{-3n}n^{3n+\frac{3}{2}}}{c3^{3n+\frac{1}{2}}n^{3n+\frac{1}{2}}e^{-3n}}\\
                                        &= c^2n3^{-(3n+\frac{1}{2})}\\
                                        &= nc^23^{-\frac{1}{2}}(3^{-3})^{n}
            \end{align*}
            so taking $n$-th roots we get
             \[
            \limsup_{n\to \infty}(a_n)^{\frac{1}{n}} \approx 3^{-3}
            \] 
            so $R = 3^{3} = 27$
    \end{enumerate}
}

\section{Exercise 1.18}
\emph{
    Let $f$ be a power series centered at the origin. Prove that $f$ has a power series expansion around any point in its disc of convergence.
}
\solution{
    Let $f(z) = \sum_{n=1}^{\infty}a_nz^{n}$ with radius of covergence $R$ then let  $z_0$ be a point in $D_R(0)$ the open disc of radius  $R$ centered at 0. we have
    \[
    f(z) = \sum_{n=0}^{\infty}a_n(z_0 + (z-z_0))^{n} = \sum_{n=0}^{\infty}a_n\sum_{k=0}^{n}\binom{n}{k}(z_0)^{n-k}(z-z_0)^{k}
    \] 
    We can swap the order of summation because for $|z| < R$ the inside converges absolutely
    \[
    f(z) = \sum_{k=0}^{\infty}\left(\sum_{n=k}^{\infty}\binom{n}{k}(z_0)^{n-k}a_n\right)(z-z_0)^{k}
    \] 
    We know that the original series converges if $|z| < R$ so we need
     \[
    |z_0 - (z-z_0)| < R \implies |z-z_0| < R - |z_0|
    \] 
    so the radius of convergence is at least $R - |z_0|$
}

\section{Exercise 1.20}
\emph{
    Expand $(1-z)^{-m}$ in powers of $z$. Here $m$ is a fixed positive integer. Also show that if
    \[
        (1-z)^{-m} = \sum_{n=0}^{\infty}a_nz^{n}
    \] 
    then one obtains the following asymptotic relation for the coefficients:
    \[
    a_n \sim \frac{1}{(m-1)!}n^{m-1} \qquad n \rightarrow \infty
    \] 
}
\solution{
    $f(0) = 1$,  $f'(z) = m(1-z)^{-m-1}$ $f'(0) = m$  $f^{(2)}(z) = m(m+1)(1-z)^{-(m+2)}$ 
    in general
    \[
    f^{(n)}(z) = \frac{(m+n-1)!}{(m-1)!}(1-z)^{-(m+n)} \qquad f^{(n)}(0) = \frac{(m+n-1)!}{(m-1)!}
    \] 
    so we get
    \[
    f(z) = \sum_{n=0}^{\infty}\frac{(m+n-1)!}{(m-1)!n!}z^{n}
    \] 
    we want to show that
    \[
        \frac{(m+n-1)!}{(m-1)!n!} \sim \frac{1}{(m-1)!}n^{m-1}
    \] 
    this is true if 
    \[
        \frac{(m+n-1)!}{n!} \sim n^{m-1}
    \] 
    but we have
    \[
        \frac{(m+n-1)!}{n!} = (m+n-1)(m+n-2)\dots(n+1) = n^{m-1} + O(n^{-m_2})
    \] 
    so in the limit they approach eachother
}

\section{Exercise 1.21}
\emph{
    Show that for $|z| < 1$ one has
     \[
    \frac{z}{1-z^2} + \frac{z^2}{1-z^{4}} + \dots + \frac{z^{2^{n}}}{1-z^{2^{n+1}}} + \dots = \frac{z}{1-z}
    \] 
    and
    \[
    \frac{z}{1+z} + \frac{2z^2}{1+z^2} + \dots + \frac{2^{k}z^{2^{k}}}{1+z^{2^{k}}} + \dots = \frac{z}{1-z}
    \] 
    Justify any change in the order of summation
}
\solution{
    For the first sum we have
    \[
    \sum_{n=0}^{\infty}\frac{z^{2^{n}}}{1-z^{2^{n+1}}}
    \] 
    We try to write this as a telescoping sum
    \[
        \frac{z^{2^{n}}}{1-z^{2^{n+1}}} = x - \frac{1}{1-z^{2^{n+1}}}
    \] 
    \[
    z^{2^{n}} = (1-z^{2^{n+1}})x - 1
    \] 
    \[
    x = \frac{1+z^{2^{n}}}{1-z^{2^{n+1}}} = \frac{1+z^{2^{n}}}{(1-z^{2^{n}})(1+z^{2^{n}})} = \frac{1}{1-z^{2^{n}}}
    \] 
    so our original sum is
    \[
    \sum_{n=0}^{\infty}\left(\frac{1}{1-z^{2^{n}}} - \frac{1}{1-z^{2^{n+1}}}\right)
    \] 
    so the m-th partial sum is
    \[
        \frac{1}{1-z} - \frac{1}{1-z^{2^{m+1}}}
    \] 
    as $m \rightarrow \infty$ this converges to
    \[
    \frac{1}{1-z} - 1 = \frac{z}{1-z}
    \] 
    our desired result.
    For the second series we similarly want to find $x_k$ such that
     \[
         x_k - x_{k+1} = \frac{2^{k}z^{2^{k}}}{1+z^{2^{k}}}
    \] 
    we try $x_k = \frac{2^{k}z^{2^{k}}}{1-z^{2^{k}}}$
    then
    \begin{align*}
        x_k - x_{k+1} &= \frac{2^{k}z^{2^{k}}}{1-z^{2^{k}}} - \frac{2^{k+1}z^{2^{k+1}}}{1-z^{2^{k+1}}}\\
                      &= \frac{2^{k}z^{2^{k}}}{1-z^{2^{k}}} - \frac{2^{k+1}z^{2^{k+1}}}{(1-z^{2^{k}})(1+z^{2^{k}})}\\
                      &= \frac{2^{k}z^{2^{k}}+2^{k}z^{2^{k+1}} - 2^{k+1}z^{2^{k+1}}}{1-z^{2^{k+1}}}\\
                      &= 2^{k}\frac{z^{2^{k}} - z^{2^{k+1}}}{(1-z^{2^{k}})(1+z^{2^{k}})}\\
                      &= 2^{k}z^{2^{k}}\frac{(1-z^{2^{k})}}{(1-z^{2^{k}})(1+z^{2^{k}})}\\
                      &= \frac{2^{k}z^{2^{k}}}{1+z^{2^{k}}}
    \end{align*}
    as desired thus by telescoping, our  $m$-th term is
    \[
    \frac{z}{1-z} - \frac{2^{m+1}z^{2^{m+1}}}{1-z^{2^{m+1}}}
    \] 
    the ending term converges to 0. since the bottom of the fraction converges to 1 and the top goes to zero since $|z| < 1$  and $z$ to the power of an exponential shrinks faster than the exponential growth of  $2^{m}$
    thus the limit is
    \[
    \frac{z}{1-z}
    \] 
    as required.
}
\section{Exercise 23}
\emph{
    Consider the function $f$ defined on $ \mathbb{R}$ by
    \[
    f(x) = 
    \begin{cases}
        0 & x \le 0\\
        e^{-\frac{1}{x^2}} & x >0
    \end{cases}
    \] 
    Prove that $f$ is indefinitely differentiable on $ \mathbb{R}$ and that $f^{(n)}(0) = 0$ for all $n \ge 1$. Conclude
    that  $f$ does not have a converging power series expansion for $x$ near the origin.
}
\solution{
    It suffices to check for differentiability at the origin since for $x < 0$ all derivatives are 0 
    and for  $x > 0$ the function  $e^{-\frac{1}{x^2}}$ is smooth and indefinitely differentiable.
    Additionally its derivative is always a linear combination of $e^{-\frac{1}{x^2}}$ terms divided by polynomials. and so
    \[
    \lim_{x\to 0^{+}}\frac{f^{(n)}(x)}{x} = \lim_{x\to 0^{+}}f^{(n)} = 0
    \] 
    and similarly
    \[
    \lim_{x\to 0^{-}}\frac{f^{(n)}(x)}{x} = \lim_{x\to 0^{-}}f^{(n)} = 0
    \] 
    and so
    \[
    f^{(n)}(0) = \lim_{h\to 0} \frac{f^{(n-1)}(-h) + f^{(n)}(h)}{h} = 0
    \] 
    for all $n$.
    Assume for contradiction there exists an expansion that converges in a radius $ r >0$ around the origin
     \[
    \sum_{n=0}^{\infty}a_nx^{n}
    \] 
    then by theorem $2.6$ since  $f$ is also a power series on $ \mathbb{C}$ we have that the derivates
    \[
    f^{(m)}(x) = \sum_{n=m}^{\infty}n(n-1)\dots(n-m+1)a_nz^{n-m}
    \] 
    however by plugging in $0$ we get that $a_n = 0$ for all  $n$ and so  $f = 0$ in the ball around  $0$ however  $e^{-\frac{1}{x^2}}$ is strictly greater than $0$ for all  $x > 0$ and which is a contradiction. Thus no power series is possible
}

\end{document}
